% ==============================================================================
% WEEK 1: DIFFERENTIAL OPERATORS
% ==============================================================================

\chapter{Week 1 - Differential operators}

\section{Review, Notations, and Terminology}

We begin by recalling standard foundational sets and indexing conventions:
\begin{itemize}
    \item $\mathbb{N} = \{0, 1, 2, 3, \dots\}$ denotes the set of natural numbers.
    \item $\mathbb{Z} = \{\dots, -2, -1, 0, 1, 2, \dots\}$ denotes the set of integers.
    \item $\mathbb{R}$ denotes the field of real numbers.
\end{itemize}
For multidimensional spaces $\mathbb{R}^n = \underbrace{\mathbb{R} \times \dots \times \mathbb{R}}_{n \text{ times}}$, an arbitrary point $x \in \mathbb{R}^n$ is written in Cartesian coordinates as:
\begin{align}
    n &= 2 : \quad (x_1, x_2) = (x, y) \in \mathbb{R}^2, \\
    n &= 3 : \quad (x_1, x_2, x_3) = (x, y, z) \in \mathbb{R}^3, \\
    n > 0 : \quad &(x_1, x_2, \dots, x_n) \in \mathbb{R}^n.
\end{align}

\begin{definition}{Scalar Field}{scalar_field}
Let $\Omega \subset \mathbb{R}^n$ be an open domain. A real-valued function
\begin{equation}
    f : \Omega \longrightarrow \mathbb{R}, \quad x \longmapsto f(x) = f(x_1, x_2, \dots, x_n)
\end{equation}
depending on multiple variables is called a \textbf{scalar field} defined on $\Omega$.
\end{definition}

\begin{remarque}[Continuity and Differentiability Classes]
For $k \in \mathbb{N}$, we write $f \in \mathcal{C}^k(\Omega)$ if all partial derivatives of order $\le k$ exist and are continuous on $\Omega$.
\begin{itemize}
    \item For $k = 0$, $f \in \mathcal{C}^0(\Omega)$ signifies that $f$ is continuous on $\Omega$.
    \item For instance, $f(x) = \sqrt{x}$ on $(0, \infty)$ satisfies $f'(x) = \frac{1}{2\sqrt{x}}$, so $f \in \mathcal{C}^1\bigl((0, \infty)\bigr)$.
\end{itemize}
\end{remarque}

\begin{definition}{Vector Field}{vector_field}
Let $\Omega \subset \mathbb{R}^n$ be an open domain with $n \ge 1$ and $m > 1$. A vector-valued function
\begin{equation}
    F : \Omega \longrightarrow \mathbb{R}^m, \quad x \longmapsto F(x) = \big(F_1(x), F_2(x), \dots, F_m(x)\big)
\end{equation}
where each component $F_i : \Omega \to \mathbb{R}$ is a scalar field, is called a \textbf{vector field}. A vector field can thus be viewed as an ordered collection of scalar fields. We write $F \in \mathcal{C}^k(\Omega, \mathbb{R}^m)$ if every component satisfies $F_i \in \mathcal{C}^k(\Omega)$ for all $i = 1, \dots, m$.
\end{definition}

\begin{definition}{Nabla Operator (Del)}{nabla_operator}
The formal vector differential operator \textbf{nabla} (or \textbf{del}), denoted by $\vec{\nabla}$ or $\nabla$, is defined in $\mathbb{R}^n$ by:
\begin{equation}
    \vec{\nabla} = \left( \pdv{}{x_1}, \, \pdv{}{x_2}, \, \dots, \, \pdv{}{x_n} \right).
\end{equation}
\end{definition}

---

\section{The Gradient Operator}

\begin{definition}{Gradient of a Scalar Field}{gradient_def}
Let $\Omega \subset \mathbb{R}^n$ be an open domain and let $f : \Omega \to \mathbb{R}$ be a scalar field such that $f \in \mathcal{C}^1(\Omega)$. The \textbf{gradient} of $f$, denoted by $\mathrm{grad}\,f$, $\vec{\nabla} f$, or $\mathrm{D}f$, is the vector field defined by:
\begin{equation}
    \mathrm{grad}\,f(x) = \vec{\nabla} f(x) = \left( \pdv{f}{x_1}(x), \, \pdv{f}{x_2}(x), \, \dots, \, \pdv{f}{x_n}(x) \right) \in \mathbb{R}^n.
\end{equation}
Formally, it represents the symbolic product of the nabla operator $\vec{\nabla}$ with the scalar function $f$:
\begin{equation}
    \vec{\nabla} f = \left( \pdv{}{x_1}, \, \dots, \, \pdv{}{x_n} \right) \cdot f = \left( \pdv{f}{x_1}, \, \dots, \, \pdv{f}{x_n} \right).
\end{equation}
The gradient transforms a scalar field into a vector field: $\mathrm{grad} : \mathcal{C}^1(\Omega, \mathbb{R}) \to \mathcal{C}^0(\Omega, \mathbb{R}^n)$.
\end{definition}

\begin{remarque}[Three-Dimensional Gradient]
In $\mathbb{R}^3$, where $(x_1, x_2, x_3) = (x, y, z)$, the gradient evaluates to:
\begin{equation}
    \vec{\nabla} f(x, y, z) = \left( \pdv{f}{x}, \, \pdv{f}{y}, \, \pdv{f}{z} \right).
\end{equation}
\end{remarque}

\begin{exemple}[Direct Computation of Gradient]
Consider $f : \mathbb{R}^3 \to \mathbb{R}$ defined by $f(x, y, z) = x^2 y^3 \sin(z^2)$. The partial derivatives are:
\begin{align}
    \pdv{f}{x} &= 2x y^3 \sin(z^2), \\
    \pdv{f}{y} &= 3x^2 y^2 \sin(z^2), \\
    \pdv{f}{z} &= 2x^2 y^3 z \cos(z^2).
\end{align}
Hence:
\begin{equation}
    \vec{\nabla} f(x, y, z) = \Big( 2x y^3 \sin(z^2), \, 3x^2 y^2 \sin(z^2), \, 2x^2 y^3 z \cos(z^2) \Big).
\end{equation}
\end{exemple}

\begin{exemple}[Gravitational Potential Field]
Let $P_0 = (x_0, y_0, z_0)$ be a source point bearing a mass $M$, and $P = (x, y, z)$ the observation point bearing a mass $m$. The distance between them is:
\begin{equation}
    r = \norm{P - P_0} = \sqrt{(x - x_0)^2 + (y - y_0)^2 + (z - z_0)^2}.
\end{equation}
The gravitational potential $f : \mathbb{R}^3 \setminus \{P_0\} \to \mathbb{R}$ is given by:
\begin{equation}
    f(x, y, z) = \psi(r) = \frac{C}{r}, \quad \text{where } C = G M m.
\end{equation}
Applying the chain rule to compute the partial derivatives:
\begin{equation}
    \pdv{r}{x} = \pdv{}{x} \sqrt{(x - x_0)^2 + (y - y_0)^2 + (z - z_0)^2} = \frac{x - x_0}{r}.
\end{equation}
Since $\frac{\dd\psi}{\dd r} = -\frac{C}{r^2}$, we obtain:
\begin{equation}
    \pdv{f}{x} = \frac{\dd\psi}{\dd r} \pdv{r}{x} = -\frac{C}{r^2} \cdot \frac{x - x_0}{r} = -\frac{C(x - x_0)}{r^3}.
\end{equation}
By radial symmetry:
\begin{equation}
    \pdv{f}{y} = -\frac{C(y - y_0)}{r^3}, \qquad \pdv{f}{z} = -\frac{C(z - z_0)}{r^3}.
\end{equation}
The resulting gravitational force field is the gradient of the potential:
\begin{equation}
    \vec{\nabla} f(x, y, z) = -\frac{C}{r^3} \big( x - x_0, \, y - y_0, \, z - z_0 \big) = -\frac{C}{r^3} (P - P_0).
\end{equation}
\end{exemple}

---

\section{The Divergence Operator}

\begin{definition}{Divergence of a Vector Field}{divergence_def}
Let $\Omega \subset \mathbb{R}^n$ be an open domain and let $F : \Omega \to \mathbb{R}^n$ be a vector field with $F \in \mathcal{C}^1(\Omega, \mathbb{R}^n)$. The \textbf{divergence} of $F$, denoted by $\mathrm{div}\,F$ or $\vec{\nabla} \cdot F$, is the scalar field defined by:
\begin{equation}
    \mathrm{div}\,F(x) = \vec{\nabla} \cdot F(x) = \sum_{i=1}^n \pdv{F_i}{x_i}(x) = \pdv{F_1}{x_1} + \pdv{F_2}{x_2} + \dots + \pdv{F_n}{x_n} \in \mathbb{R}.
\end{equation}
The divergence operator maps vector fields to scalar fields: $\mathrm{div} : \mathcal{C}^1(\Omega, \mathbb{R}^n) \to \mathcal{C}^0(\Omega, \mathbb{R})$.
\end{definition}

\begin{remarque}[Physical Interpretation: Flux and Conservation]
Consider an infinitesimal control volume in 2D of dimensions $\dd x \times \dd y$ centered around $(x, y)$ with flow velocity $V = (V_x, V_y)$:
\begin{itemize}
    \item Net outflow along the $x$-direction:
    \begin{equation}
        \left( V_x + \pdv{V_x}{x} \dd x \right) \dd y - V_x \dd y = \pdv{V_x}{x} \dd x \dd y.
    \end{equation}
    \item Net outflow along the $y$-direction:
    \begin{equation}
        \left( V_y + \pdv{V_y}{y} \dd y \right) \dd x - V_y \dd x = \pdv{V_y}{y} \dd x \dd y.
    \end{equation}
\end{itemize}
Summing both components yields the total rate of outward mass expansion per unit area:
\begin{equation}
    \text{Net Rate of Change} = \left( \pdv{V_x}{x} + \pdv{V_y}{y} \right) \dd x \dd y = (\vec{\nabla} \cdot V) \dd x \dd y.
\end{equation}
Consequently:
\begin{itemize}
    \item $\vec{\nabla} \cdot V(x) > 0$: The point acts as a \textbf{source} (net outward expansion).
    \item $\vec{\nabla} \cdot V(x) < 0$: The point acts as a \textbf{sink} (net inward convergence).
    \item $\vec{\nabla} \cdot V(x) = 0$: The fluid field is \textbf{incompressible} (solenoidal).
\end{itemize}
\end{remarque}

\begin{exemple}[Direct Computation of Divergence]
Let $F : \mathbb{R}^3 \to \mathbb{R}^3$ be given by $F(x, y, z) = \big(x^2 - e^y, \, \sin(xz), \, y^2 e^{2xz}\big)$. Then:
\begin{align}
    \mathrm{div}\,F(x, y, z) &= \pdv{}{x}\left(x^2 - e^y\right) + \pdv{}{y}\big(\sin(xz)\big) + \pdv{}{z}\left(y^2 e^{2xz}\right) \nonumber \\
    &= 2x + 0 + 2x y^2 e^{2xz} \nonumber \\
    &= 2x\left(1 + y^2 e^{2xz}\right).
\end{align}
\end{exemple}

\begin{exemple}[Divergence of an Inverse-Square Central Force Field]
Let $P_0 = (0, 0, 0)$, $r = \sqrt{x^2 + y^2 + z^2}$, and consider the central force field $F : \mathbb{R}^3 \setminus \{(0,0,0)\} \to \mathbb{R}^3$:
\begin{equation}
    F(x, y, z) = -\frac{C}{r^3} (x, y, z) = \left( -\frac{Cx}{r^3}, \, -\frac{Cy}{r^3}, \, -\frac{Cz}{r^3} \right).
\end{equation}
Differentiating the first component $F_1 = -C x r^{-3}$ with respect to $x$:
\begin{align}
    \pdv{F_1}{x} &= -C \left( \frac{1}{r^3} - \frac{3x}{r^4} \pdv{r}{x} \right) = -C \left( \frac{1}{r^3} - \frac{3x^2}{r^5} \right) = \frac{C(3x^2 - r^2)}{r^5}.
\end{align}
By identical permutation for $y$ and $z$:
\begin{equation}
    \pdv{F_2}{y} = \frac{C(3y^2 - r^2)}{r^5}, \qquad \pdv{F_3}{z} = \frac{C(3z^2 - r^2)}{r^5}.
\end{equation}
Summing the three partial derivatives:
\begin{align}
    \mathrm{div}\,F &= \frac{C}{r^5} \Big( (3x^2 - r^2) + (3y^2 - r^2) + (3z^2 - r^2) \Big) \nonumber \\
    &= \frac{C}{r^5} \Big( 3(x^2 + y^2 + z^2) - 3r^2 \Big) = \frac{C}{r^5}(3r^2 - 3r^2) = 0.
\end{align}
Thus, the inverse-square gravitational field is divergence-free on its entire domain of definition $\mathbb{R}^3 \setminus \{(0,0,0)\}$.
\end{exemple}

---

\section{The Curl Operator (Rotor)}

\begin{definition}{Curl Operator}{curl_def}
Let $\Omega \subset \mathbb{R}^n$ be an open domain and let $F : \Omega \to \mathbb{R}^n$ be a vector field of class $\mathcal{C}^1(\Omega, \mathbb{R}^n)$. The \textbf{curl} operator (also denoted $\mathrm{rot}\,F$ or $\vec{\nabla} \times F$) is defined as follows:

\begin{itemize}
    \item \textbf{In 2D} ($n = 2$), for $F = (F_1, F_2) : \Omega \subset \mathbb{R}^2 \to \mathbb{R}^2$, the scalar curl is:
    \begin{equation}
        \mathrm{curl}\,F(x, y) = \pdv{F_2}{x}(x, y) - \pdv{F_1}{y}(x, y) \in \mathbb{R}.
    \end{equation}
    \item \textbf{In 3D} ($n = 3$), for $F = (F_1, F_2, F_3) : \Omega \subset \mathbb{R}^3 \to \mathbb{R}^3$, the curl is a vector field:
    \begin{equation}
        \mathrm{curl}\,F(x, y, z) = \left( \pdv{F_3}{y} - \pdv{F_2}{z}, \, \pdv{F_1}{z} - \pdv{F_3}{x}, \, \pdv{F_2}{x} - \pdv{F_1}{y} \right) \in \mathbb{R}^3.
    \end{equation}
\end{itemize}
\end{definition}

\begin{remarque}[Formal Determinant Representation]
In $\mathbb{R}^3$, the curl can be computed via the formal cross product with nabla using a symbolic $3 \times 3$ determinant:
\begin{equation}
    \vec{\nabla} \times F = 
    \begin{vmatrix}
        \bm{e}_x & \bm{e}_y & \bm{e}_z \\
        \pdv{}{x} & \pdv{}{y} & \pdv{}{z} \\
        F_1 & F_2 & F_3
    \end{vmatrix}
    = \left( \pdv{F_3}{y} - \pdv{F_2}{z} \right) \bm{e}_x + \left( \pdv{F_1}{z} - \pdv{F_3}{x} \right) \bm{e}_y + \left( \pdv{F_2}{x} - \pdv{F_1}{y} \right) \bm{e}_z.
\end{equation}
Furthermore, if $F$ is strictly two-dimensional embedded in 3D, such that $F_3 = 0$ and $\pdv{F_1}{z} = \pdv{F_2}{z} = 0$, then:
\begin{equation}
    \vec{\nabla} \times F = \Big( 0, \, 0, \, \mathrm{curl}_{2D}(F_1, F_2) \Big).
\end{equation}
\end{remarque}

\begin{exemple}[Curl of an Inverse-Square Central Force Field]
Let $F(x, y, z) = -\frac{C}{r^3}(x, y, z)$ on $\mathbb{R}^3 \setminus \{(0,0,0)\}$, where $r = \sqrt{x^2+y^2+z^2}$. Computing the first component of $\vec{\nabla} \times F$:
\begin{align}
    (\vec{\nabla} \times F)_1 &= \pdv{}{y}\left(-\frac{C z}{r^3}\right) - \pdv{}{z}\left(-\frac{C y}{r^3}\right) \nonumber \\
    &= -C z \left(-\frac{3}{r^4} \pdv{r}{y}\right) - \left[ -C y \left(-\frac{3}{r^4} \pdv{r}{z}\right) \right] \nonumber \\
    &= \frac{3C z y}{r^5} - \frac{3C y z}{r^5} = 0.
\end{align}
By identical cyclic permutations:
\begin{equation}
    (\vec{\nabla} \times F)_2 = \frac{3C(x z - z x)}{r^5} = 0, \qquad (\vec{\nabla} \times F)_3 = \frac{3C(y x - x y)}{r^5} = 0.
\end{equation}
Hence:
\begin{equation}
    \vec{\nabla} \times F = \begin{pmatrix} 0 \\ 0 \\ 0 \end{pmatrix}.
\end{equation}
An inverse-square central force field is strictly \textbf{irrotational}.
\end{exemple}

---

\section{The Laplacian Operator}

\begin{definition}{Laplacian of a Scalar Field}{laplacian_def}
Let $\Omega \subset \mathbb{R}^n$ be an open domain and let $f : \Omega \to \mathbb{R}$ be a scalar field such that $f \in \mathcal{C}^2(\Omega)$. The \textbf{Laplacian} operator of $f$, denoted by $\Delta f$ or $\nabla^2 f$, is defined by:
\begin{equation}
    \Delta f(x) = \sum_{i=1}^n \pdv[2]{f}{x_i}(x) = \pdv[2]{f}{x_1}(x) + \pdv[2]{f}{x_2}(x) + \dots + \pdv[2]{f}{x_n}(x) \in \mathbb{R}.
\end{equation}
The Laplacian can be viewed as the divergence of the gradient:
\begin{equation}
    \Delta f = \vec{\nabla} \cdot (\vec{\nabla} f) = \mathrm{div}(\mathrm{grad}\,f).
\end{equation}
\end{definition}

\begin{exemple}[Laplacian of the Gravitational Potential]
Let $f : \mathbb{R}^3 \setminus \{(0,0,0)\} \to \mathbb{R}$ be the gravitational potential $f(x, y, z) = \frac{C}{r}$, with $r = \sqrt{x^2 + y^2 + z^2}$. First, evaluate the first-order partial derivative:
\begin{equation}
    \pdv{f}{x} = \frac{\dd}{\dd r}\left(\frac{C}{r}\right) \pdv{r}{x} = -\frac{C}{r^2} \cdot \frac{x}{r} = -\frac{C x}{r^3}.
\end{equation}
Computing the second derivative with respect to $x$:
\begin{align}
    \pdv[2]{f}{x} &= \pdv{}{x}\left(-\frac{C x}{r^3}\right) = -C \left( \frac{1}{r^3} - \frac{3x}{r^4} \pdv{r}{x} \right) = \frac{C(3x^2 - r^2)}{r^5}.
\end{align}
Summing the second derivatives along all three Cartesian directions:
\begin{align}
    \Delta f &= \pdv[2]{f}{x} + \pdv[2]{f}{y} + \pdv[2]{f}{z} \nonumber \\
    &= \frac{C}{r^5} \Big( (3x^2 - r^2) + (3y^2 - r^2) + (3z^2 - r^2) \Big) \nonumber \\
    &= \frac{C}{r^5} \Big( 3(x^2 + y^2 + z^2) - 3r^2 \Big) = 0.
\end{align}
Therefore, outside the origin, the gravitational potential satisfies Laplace's equation $\Delta f = 0$ and is a \textbf{harmonic function}.
\end{exemple}

\section{Special Operator Combinations and Proofs}

\subsection{Divergence of a Scaled Gradient}
\begin{proposition}{Divergence of $f \grad g$}{div_f_grad_g}
Let $f \in \mathcal{C}^1(U, \R)$ and $g \in \mathcal{C}^2(U, \R)$. Then:
\begin{equation}
    \dive (f \grad g) = f \laplacian g + \inner{\grad f}{\grad g}.
\end{equation}
\end{proposition}

\begin{demonstration}
In index notation, with the summation convention implied over $i \in \{1, 2, 3\}$:
\begin{align*}
    \dive (f \grad g) &= \partial_i \left( f \, \partial_i g \right) \\
    &= (\partial_i f)(\partial_i g) + f (\partial_i \partial_i g) \\
    &= \inner{\grad f}{\grad g} + f \laplacian g. \qedhere
\end{align*}
\end{demonstration}

\subsection{Gradient of a Scalar Product}
\begin{proposition}{Gradient of $fg$}{grad_fg}
Let $f, g \in \mathcal{C}^1(U, \R)$. Then:
\begin{equation}
    \grad (fg) = f \grad g + g \grad f.
\end{equation}
\end{proposition}

\begin{demonstration}
Evaluating component-wise for each $i \in \{1, 2, 3\}$:
\begin{equation*}
    [\grad(fg)]_i = \partial_i (fg) = (\partial_i f) g + f (\partial_i g) = [g \grad f + f \grad g]_i.
\end{equation*}
Thus:
\begin{equation*}
    \grad(fg) = f \grad g + g \grad f. \qedhere
\end{equation*}
\end{demonstration}

\subsection{Divergence of a Scaled Vector Field}
\begin{proposition}{Divergence of $f \bm{F}$}{div_f_F}
Let $f \in \mathcal{C}^1(U, \R)$ and $\bm{F} \in \mathcal{C}^1(U, \R^3)$. Then:
\begin{equation}
    \dive (f \bm{F}) = f \dive \bm{F} + \inner{\grad f}{\bm{F}}.
\end{equation}
\end{proposition}

\begin{demonstration}
Using summation over repeated indices $i \in \{1, 2, 3\}$:
\begin{align*}
    \dive (f \bm{F}) &= \partial_i (f F_i) \\
    &= (\partial_i f) F_i + f (\partial_i F_i) \\
    &= \inner{\grad f}{\bm{F}} + f \dive \bm{F}. \qedhere
\end{align*}
\end{demonstration}

\subsection{Curl of a Curl}
\begin{proposition}{Curl of Curl Identity}{rot_rot_F}
Let $\bm{F} \in \mathcal{C}^2(U, \R^3)$. Then:
\begin{equation}
    \rot (\rot \bm{F}) = -\laplacian \bm{F} + \grad(\dive \bm{F}).
\end{equation}
\end{proposition}

\begin{demonstration}
Using the Levi-Civita permutation symbol $\varepsilon_{ijk}$, where $[\rot \bm{F}]_j = \varepsilon_{jlm} \partial_l F_m$:
\begin{align*}
    [\rot(\rot \bm{F})]_i &= \varepsilon_{ijk} \partial_j (\varepsilon_{klm} \partial_l F_m) \\
    &= \varepsilon_{kij} \varepsilon_{klm} \partial_j \partial_l F_m.
\end{align*}
Applying the contraction identity $\varepsilon_{kij} \varepsilon_{klm} = \delta_{il}\delta_{jm} - \delta_{im}\delta_{jl}$:
\begin{align*}
    [\rot(\rot \bm{F})]_i &= (\delta_{il}\delta_{jm} - \delta_{im}\delta_{jl}) \partial_j \partial_l F_m \\
    &= \partial_j \partial_i F_j - \partial_j \partial_j F_i.
\end{align*}
By Schwarz's theorem for $\mathcal{C}^2$ fields, $\partial_j \partial_i F_j = \partial_i (\partial_j F_j) = [\grad(\dive \bm{F})]_i$, and $\partial_j \partial_j F_i = [\laplacian \bm{F}]_i$:
\begin{equation*}
    \rot(\rot \bm{F}) = \grad(\dive \bm{F}) - \laplacian \bm{F} = -\laplacian \bm{F} + \grad(\dive \bm{F}). \qedhere
\end{equation*}
\end{demonstration}

\subsection{Curl of a Scaled Vector Field}
\begin{proposition}{Curl of $f \bm{F}$}{rot_f_F}
Let $f \in \mathcal{C}^1(U, \R)$ and $\bm{F} \in \mathcal{C}^1(U, \R^3)$. Then:
\begin{equation}
    \rot (f \bm{F}) = \grad f \times \bm{F} + f \rot \bm{F}.
\end{equation}
\end{proposition}

\begin{demonstration}
Using the explicit cross-product formulation:
\begin{equation*}
    \rot (f \bm{F}) = \begin{pmatrix} 
        \partial_y(f F_z) - \partial_z(f F_y) \\ 
        \partial_z(f F_x) - \partial_x(f F_z) \\ 
        \partial_x(f F_y) - \partial_y(f F_x) 
    \end{pmatrix}.
\end{equation*}
Applying the product rule to the $x$-component:
\begin{align*}
    \partial_y(f F_z) - \partial_z(f F_y) &= (\partial_y f) F_z + f (\partial_y F_z) - (\partial_z f) F_y - f (\partial_z F_y) \\
    &= \bigl[(\grad f) \times \bm{F}\bigr]_x + f \bigl[\rot \bm{F}\bigr]_x.
\end{align*}
Repeating identically for the $y$- and $z$-components yields:
\begin{equation*}
    \rot(f \bm{F}) = \grad f \times \bm{F} + f \rot \bm{F}. \qedhere
\end{equation*}
\end{demonstration}

\section{The Laplacian in Polar Coordinates}

\subsection{Derivation via the Multivariate Chain Rule}
\begin{theoreme}{Polar Coordinate Representation of $\laplacian f$}{laplacian_polar}
Let $x = r \cos\theta$ and $y = r \sin\theta$ define the standard transformation from polar coordinates $(r, \theta) \in (0, \infty) \times [0, 2\pi)$ to Cartesian coordinates $(x, y) \in \R^2 \setminus \{(0,0)\}$. For a scalar field $f \in \mathcal{C}^2(\R^2, \R)$, let $g(r, \theta) = f(r \cos\theta, r \sin\theta)$. Then the Laplace operator satisfies:
\begin{equation}
    \laplacian f = \pdv[2]{g}{r} + \frac{1}{r}\pdv{g}{r} + \frac{1}{r^2}\pdv[2]{g}{\theta}.
\end{equation}
\end{theoreme}

\begin{demonstration}
Applying the first-order multivariate chain rule to $g(r, \theta) = f(x(r,\theta), y(r,\theta))$:
\begin{align}
    \pdv{g}{r} &= \cos\theta \pdv{f}{x} + \sin\theta \pdv{f}{y}, \\
    \pdv{g}{\theta} &= -r\sin\theta \pdv{f}{x} + r\cos\theta \pdv{f}{y}.
\end{align}
Computing the second partial derivative with respect to $r$:
\begin{align}
    \pdv[2]{g}{r} &= \cos\theta \pdv{}{r}\left(\pdv{f}{x}\right) + \sin\theta \pdv{}{r}\left(\pdv{f}{y}\right) \notag \\
    &= \cos\theta \left( \cos\theta \pdv[2]{f}{x} + \sin\theta \pdv[2]{f}{y}{x} \right) + \sin\theta \left( \cos\theta \pdv[2]{f}{x}{y} + \sin\theta \pdv[2]{f}{y} \right) \notag \\
    &= \cos^2\theta \pdv[2]{f}{x} + 2\sin\theta\cos\theta \pdv[2]{f}{x}{y} + \sin^2\theta \pdv[2]{f}{y}.
\end{align}
Computing the second partial derivative with respect to $\theta$:
\begin{align}
    \pdv[2]{g}{\theta} &= \pdv{}{\theta}\left(-r\sin\theta \pdv{f}{x} + r\cos\theta \pdv{f}{y}\right) \notag \\
    &= -r\cos\theta \pdv{f}{x} - r\sin\theta \pdv{}{\theta}\left(\pdv{f}{x}\right) - r\sin\theta \pdv{f}{y} + r\cos\theta \pdv{}{\theta}\left(\pdv{f}{y}\right) \notag \\
    &= -r \left(\cos\theta \pdv{f}{x} + \sin\theta \pdv{f}{y}\right) \notag \\
    &\quad + r \left[ -\sin\theta \left(-r\sin\theta \pdv[2]{f}{x} + r\cos\theta \pdv[2]{f}{y}{x}\right) + \cos\theta \left(-r\sin\theta \pdv[2]{f}{x}{y} + r\cos\theta \pdv[2]{f}{y}\right) \right] \notag \\
    &= -r \pdv{g}{r} + r^2 \left( \sin^2\theta \pdv[2]{f}{x} - 2\sin\theta\cos\theta \pdv[2]{f}{x}{y} + \cos^2\theta \pdv[2]{f}{y} \right).
\end{align}
Dividing equation (2.16) by $r^2$:
\begin{equation}
    \frac{1}{r^2}\pdv[2]{g}{\theta} = -\frac{1}{r}\pdv{g}{r} + \sin^2\theta \pdv[2]{f}{x} - 2\sin\theta\cos\theta \pdv[2]{f}{x}{y} + \cos^2\theta \pdv[2]{f}{y}.
\end{equation}
Summing equations (2.15) and (2.17):
\begin{align*}
    \pdv[2]{g}{r} + \frac{1}{r}\pdv{g}{r} + \frac{1}{r^2}\pdv[2]{g}{\theta} &= (\cos^2\theta + \sin^2\theta)\pdv[2]{f}{x} + (\sin^2\theta + \cos^2\theta)\pdv[2]{f}{y} \\
    &= \pdv[2]{f}{x} + \pdv[2]{f}{y} \\
    &= \laplacian f. \qedhere
\end{align*}
\end{demonstration}